\section{States Studied}
\label{sec:states}

Table~\ref{tab:states} summarises the six states, their 2020 seat counts,
and the ensemble sources used for comparison.

\begin{table}[h]
\centering
\caption{States studied in G.1, with delegation size and ensemble source.}
\label{tab:states}
\begin{tabular}{llcl}
\toprule
State & Abbr. & $k$ (2020) & Ensemble source \\
\midrule
North Carolina & NC & 14 & \gc{} \recom{}, 1,000 steps (this paper) \\
Wisconsin      & WI &  8 & \gc{} \recom{}, 1,000 steps (this paper) \\
Georgia        & GA & 14 & \gc{} \recom{}, 1,000 steps (this paper) \\
Pennsylvania   & PA & 17 & \gc{} \recom{}, 1,000 steps (this paper) \\
Texas          & TX & 38 & Running --- placeholder $< 5$th percentile \lit \\
California     & CA & 52 & Running --- placeholder $< 5$th percentile \lit \\
\bottomrule
\multicolumn{4}{l}{\small \lit~Pending; placeholder based on algorithm design expectations.}
\end{tabular}
\end{table}

\subsection{Methodological note: ReCom samples uniformly; AR minimises}

\gc{} \recom{} samples redistricting plans \emph{uniformly at random} from
the feasible space (population balance + contiguity constraints).  The \ar{}
algorithm, by contrast, \emph{minimises edge cut}.  Therefore the \ar{} plan
is expected to fall at the low end of the cut-fraction distribution---
corresponding to the high end of compactness.  The 0.1--0.2nd percentile
results for Wisconsin, Georgia, and Pennsylvania confirm the algorithm is
functioning as designed: it finds the most compact feasible plans, not typical
ones.

This is a feature, not a coincidence.  Ensemble comparison for a
minimum-cut algorithm should produce low cut percentiles.  If it did not,
that would indicate a failure of the optimiser.  North Carolina's 50th
percentile is the informative exception: it reflects a tightly constrained
geography where the minimum-cut solution and the ensemble median converge.

\subsection{North Carolina ($k = 14$)}

North Carolina is the canonical benchmark state for ensemble redistricting
research.  The state has 14 congressional districts ($14 = 2 \times 7$) and
has been at the center of multiple redistricting challenges.  NC's geography
(coastal plain, piedmont, mountains) produces a moderately constrained
ensemble where compact plans converge toward similar cut fractions.

\subsection{Wisconsin ($k = 8$)}

Wisconsin is a competitive state with strong urban-rural sorting.  At
$k = 8 = 2^3$, the plan space is smaller than for larger states.  The three-level
binary factorisation (state $\to$ 2 halves $\to$ 4 quarters $\to$ 8 districts)
aligns well with Wisconsin's east-west geography, enabling aggressive
edge-cut minimisation.

\subsection{Georgia ($k = 14$)}

Georgia has 14 congressional districts and strong geographic partisan sorting:
Democratic voters are concentrated in the Atlanta metropolitan area and the
Black Belt.  This geographic structure means the minimum-cut partition produces
a small number of large, compact Democratic districts in Atlanta and a larger
number of smaller Republican districts covering the rural areas---a distribution
with very low edge cut relative to the ensemble.

\subsection{Pennsylvania ($k = 17$)}

Pennsylvania's $k = 17$ is prime, which means a direct 17-way partition of the
state graph rather than a hierarchical factorisation.  Despite losing the
hierarchy advantage, the \ar{} plan achieves 0.2nd percentile compactness,
confirming that METIS edge-cut optimisation is effective even for prime-$k$
states.

\subsection{Texas ($k = 38$) and California ($k = 52$)}

\recom{} ensembles for Texas and California are currently running.  Based on
the pattern across WI, GA, and PA---and the theoretical expectation that the
minimum-cut algorithm finds near-extremal solutions---we expect both states to
fall below the 5th percentile.  Results will replace this placeholder when
the runs complete.
